% Einheit 4 von 4 – Sachaufgaben (bank/quadratische-gleichungen/e4.jsonl, zusammenbau v0.1)
%% TODO Zweigzeile Teil 1: was hier gelernt wird (Satz in Schülersprache aus dem Einheitstitel) – Einheit 4
\einheitenkopf[e4]{Einheit 4 von 4 · Sachaufgaben}
\zweigzeile{ab Kl. 9, je nach Buch bis Kl. 10 · keine P10-Aufgabe}
\setcounter{aufgabe}{20}

%% TODO Titel: Ich-kann-Satz fehlt in der Bank (Platzhalter: Kettenname) – Nr. 21
%% TODO Anweisung: ein Satz über der ersten Teilaufgabe fehlt in der Bank – Nr. 21
\begin{aufgabe}{Sachaufgaben}
\begin{teile}
\teil Zahlenrätsel „Das Quadrat einer Zahl ist $25$“: Lösungen $5$ und $-5$ – welche passen?\\ \kreuz{nur $5$} \\ \kreuz{nur $-5$} \\ \kreuz{beide}
\end{teile}
%% TODO Darstellung für schwach fehlt in der Bank (quadratische-gleichungen-e4-k1-s1-v1; Abschnitt „Für schwache Schüler“)
\swz[3]{Das Quadrat einer Zahl $x$ ist $144$: $x^2 = 144$. Welche Zahlen?}{}
%% TODO Darstellung für schwach fehlt in der Bank (quadratische-gleichungen-e4-k1-s1-v2; Abschnitt „Für schwache Schüler“)
\swz[3]{Das Quadrat einer Zahl $x$ ist $400$: $x^2 = 400$. Welche Zahlen?}{}
%% TODO Darstellung für schwach fehlt in der Bank (quadratische-gleichungen-e4-k1-s1-v3; Abschnitt „Für schwache Schüler“)
\swz[3]{Das Quadrat einer Zahl $x$ ist $2500$: $x^2 = 2500$. Welche Zahlen?}{}
%% TODO Darstellung für schwach fehlt in der Bank (quadratische-gleichungen-e4-k1-s1-v4; Abschnitt „Für schwache Schüler“)
\swz[3]{Das Quadrat einer Zahl $x$ ist $324$: $x^2 = 324$. Welche Zahlen?}{}
%% TODO Darstellung für schwach fehlt in der Bank (quadratische-gleichungen-e4-k1-s1-v5; Abschnitt „Für schwache Schüler“)
\swz[3]{Das Quadrat einer Zahl $x$ ist $625$: $x^2 = 625$. Welche Zahlen?}{}
\swfrage{Was bleibt gleich, was ändert sich?}
%% TODO Darstellung für schwach fehlt in der Bank (quadratische-gleichungen-e4-k1-s2-v1; Abschnitt „Für schwache Schüler“)
\swz[3]{Vermehrt man das Quadrat einer Zahl $x$ um $7$, erhält man $88$. Welche Zahlen?}{}
%% TODO Darstellung für schwach fehlt in der Bank (quadratische-gleichungen-e4-k1-s3-v1; Abschnitt „Für schwache Schüler“)
\swz[3]{Das Produkt einer Zahl $x$ mit ihrem Nachfolger ist $56$. Welche Zahlen?}{}
%% TODO Darstellung für schwach fehlt in der Bank (quadratische-gleichungen-e4-k1-s4-v1; Abschnitt „Für schwache Schüler“)
\swz[3]{Ein rechteckiger Garten ist $5$ m länger als breit und hat die Fläche $84\,\text{m}^2$. Skizze, Gleichung mit der Breite $x$, Lösungen?}{}
%% TODO Darstellung für schwach fehlt in der Bank (quadratische-gleichungen-e4-k1-s5-v1; Abschnitt „Für schwache Schüler“)
\swz[3]{Ein Rechteck ist $4$ cm länger als breit und hat die Fläche $96\,\text{cm}^2$. Wie breit ist es? Begründe, welche Lösung wegfällt.}{}
%% TODO Darstellung für schwach fehlt in der Bank (quadratische-gleichungen-e4-k1-s6-v1; Abschnitt „Für schwache Schüler“)
\swz[3]{Ein rechteckiger Tisch ist $40$ cm länger als breit und hat die Fläche $6000\,\text{cm}^2$. Breite und Länge? Antworte im Satz.}{}
%% TODO Darstellung für schwach fehlt in der Bank (quadratische-gleichungen-e4-k1-s7-v1; Abschnitt „Für schwache Schüler“)
\swz{Lukas sagt: Ein Rechteck, das $3$ cm länger als breit ist und $70\,\text{cm}^2$ Fläche hat, ist $7$ cm breit. Stimmt das?}{}
\swz[3]{Ein Rechteck ist $4$ m länger als breit. Verlängert man beide Seiten um $2$ m, beträgt die Fläche $96\,\text{m}^2$. Wie breit und wie lang ist das ursprüngliche Rechteck?}{}
\end{aufgabe}

%% TODO Titel: Ich-kann-Satz fehlt in der Bank (Platzhalter: Kettenname) – Nr. 22
%% TODO Anweisung: ein Satz über der ersten Teilaufgabe fehlt in der Bank – Nr. 22
\begin{aufgabe}{Quadrat mit Rand}
%% TODO Darstellung für schwach fehlt in der Bank (quadratische-gleichungen-e4-k2-s1-v1; Abschnitt „Für schwache Schüler“)
\swz[3]{Ein quadratisches Bild bekommt ringsum einen $5$ cm breiten Rahmen. Mit Rahmen ist die Fläche $900\,\text{cm}^2$. Wie lang ist die Seite des Bildes?}{}
\end{aufgabe}

%% TODO Titel: Ich-kann-Satz fehlt in der Bank (Platzhalter: Kettenname) – Nr. 23
%% TODO Anweisung: ein Satz über der ersten Teilaufgabe fehlt in der Bank – Nr. 23
\begin{aufgabe}{Sachaufgaben – Fehler finden, Begründen}
\swz[3]{Ein Rechteck ist $2$ m länger als breit und hat $35\,\text{m}^2$ Fläche. Wo ist der Fehler? \rechnung{x \cdot (x + 2) &= 35 \\ x^2 + 2x - 35 &= 0 \\ x_1 = 5, \quad x_2 &= -7} Antwort: Die Breite ist $-7$ m.}{}
\swfrage{Für die Breite eines Rechtecks ergibt die Gleichung $x_1 = 9$ und $x_2 = -13$. Begründe, warum nur eine Lösung die Breite ist.}
\end{aufgabe}

\uebersichtskasten{Gleichung aus Text: Unbekannte festlegen (x = …), Gleichung aus Formel oder Text aufstellen, ordnen, lösen, Lösungen an der Frage prüfen (Länge und Anzahl nicht negativ), Antwort mit Einheit. \\ Rechteck: Länge um 8 m größer als die Breite, Fläche 20 m²: x·(x + 8) = 20 → x² + 8x − 20 = 0 → x₁ = 2, x₂ = −10 \quad Breite 2 m, Länge 10 m; −10 m ist keine Länge. \\ Zahlenrätsel: beide Lösungen sind Zahlen – beide angeben, wenn die Frage nichts ausschließt.}
